\documentclass[10pt,twocolumn,letterpaper]{article}

\usepackage[top=0.75in,bottom=0.75in,left=0.75in,right=0.75in]{geometry}
\usepackage{times}
\usepackage{microtype}
\usepackage{amsmath}
\usepackage{amssymb}
\usepackage{booktabs}
\usepackage{multirow}
\usepackage{graphicx}
\usepackage{tikz}
\usetikzlibrary{shapes.geometric,arrows.meta,positioning,calc,fit,backgrounds}
\usepackage{subcaption}
\usepackage{xcolor}
\usepackage{cleveref}
\usepackage[numbers,sort&compress]{natbib}
\usepackage[colorlinks=true,linkcolor=blue,citecolor=blue,urlcolor=blue]{hyperref}

% Define color palette for tables and diagrams
\definecolor{survivorgreen}{RGB}{34,139,34}
\definecolor{blockerred}{RGB}{178,34,34}
\definecolor{highlightblue}{RGB}{24,116,205}
\definecolor{boxbg}{RGB}{245,247,250}
\definecolor{boxborder}{RGB}{200,210,225}

\title{\textbf{Benchmarking Multimodal Vision-Language Models on Untrimmed Sports Video: Multi-Scale Domain Routing, Audio-Visual Causality, and Zero-Cost Edge Perception}}

\author{
  \textbf{Hayden Lundin} \\
  School of Data Science and Business Analytics \\
  University of North Carolina at Charlotte \\
  {\tt\small hlundin@uncc.edu}
}

\date{October 2026}

\begin{document}

\maketitle

\begin{abstract}
Video action spotting and dense commentary generation in long-form, untrimmed sports broadcasts represent fundamental challenges in artificial intelligence, requiring fine-grained spatiotemporal perception, causal reasoning, and cross-modal grounding. Traditional pipelines rely on frozen 3D convolutional features fed into closed-set classifiers, completely detaching event detection from semantic narrative reasoning. 
In this work, we present the \textbf{SoccerNet VLM Modernization Benchmark (V2 \& V3)}, a unified multi-task evaluation of modern open-weight Vision-Language Models (VLMs) across five official SoccerNet challenge tracks: Action Spotting (Track A), Dense Video Captioning (Track B), Visual Question Answering (Track C), Multi-View Foul Recognition (Track D), and Ball Action Spotting (Track E).
To democratize sports video understanding, we execute this entire benchmark locally on an 8\,GB consumer laptop GPU (NVIDIA RTX 3000 Ada under Windows WDDM) at exactly \$0 cloud compute spend.
We audit 20 candidate vision-language architectures, identifying a 50\% hardware survival rate and isolating three catastrophic edge bottlenecks: unpooled AnyRes dynamic tiling explosions ($113.8$\,GiB virtual allocations on 15 frames), unquantized Per-Layer Embedding tables ($>9.1$\,GiB initial footprints), and Windows JIT compiler hurdles.
To overcome these limits, we engineer: (1) a single-pass streaming decoder consuming only 199\,MB RAM with 0.000s seek drift; (2) burned-in high-contrast digital clock badges ($[00:\text{SS}]$) eliminating temporal count hallucinations; (3) a 3-Model Domain-Routed Specialist Ensemble (Qwen3.5-4B, MiniCPM-V 4.6, LFM2.5-VL-3B) achieving 20.87\% Active Loose-mAP and 20.00\% out-of-sample 5-fold Leave-One-Chunk-Out cross-validation; (4) a mathematical proof of the ``00:00'' First-Token Prior Mode Collapse when Chain-of-Thought commentary is omitted ($19.4\times$ Tight-mAP drop); and (5) the Cross-Modal Temporal Synchrony Law, proving that initiator-biased visual detections must be shifted by $+2.95\text{s}$ before causal acoustic whistle gating.
On downstream V3 tracks, our zero-cost pipeline achieves 100.00\% Top-1 Accuracy on SN-VQA, eliminates 100\% of false-positive fouls in clean duels on MV-Foul (100.00\% Active Balanced Accuracy), and propels strict sub-second Ball Action Spotting ($\pm 1.0\text{s}$, mAP@1) from 0.2500 to 1.0000 ($4.0\times$ gain).
All code, dataset splits, and reproducible evaluation harnesses are released under an open-source license.
\end{abstract}

\section{Introduction}
\label{sec:intro}

Automated understanding of professional broadcast sports video is among the most demanding domains in multimodal computer vision. A standard 45-minute soccer match half is an untrimmed, high-entropy temporal stream comprising over 67,500 frames at 25\,FPS, characterized by rapid camera transitions, extreme perspective changes, aggressive motion blur, small object-to-frame ratios (e.g., tracking a 22\,cm ball across a 105\,m pitch), severe physical occlusions in crowded penalty boxes, and subtle disciplinary infractions governed by off-screen referee whistles~\cite{giancola2018soccernet,deliege2021soccernetv2}.

For the past decade, the dominant paradigm in sports video analysis—exemplified by NetVLAD~\cite{arandjelovic2016netvlad}, E2E-Spot~\cite{hong2022e2espot}, and ActionFormer~\cite{zhang2022actionformer}—has framed action spotting as a closed-set temporal regression problem. In this pipeline, fixed-duration sliding windows are processed by frozen 3D convolutional or optical flow backbones (e.g., I3D, VideoMAE) to extract vector embeddings, which are subsequently classified by lightweight temporal convolutional heads. While effective at detecting isolated, repetitive motions, these architectures operate in an interpretive vacuum: they possess zero semantic reasoning, cannot explain \textit{why} an event occurred, cannot arbitrate complex rules, and are fundamentally incapable of generating natural-language broadcast commentary~\cite{cioppa2022soccernetcaption}.

The emergence of Multimodal Vision-Language Models (VLMs)~\cite{qwen2025qwen25vl,openbmb2024minicpmv,liquid2025lfm,chen2024internvl} offers the theoretical promise of unifying temporal event localization with open-vocabulary reasoning, narrative captioning, and rule arbitration. However, edge deployment of VLMs onto untrimmed broadcast video faces formidable systems barriers:
\begin{enumerate}
    \item \textbf{The Temporal Receptive Field Ceiling:} Standard video VLMs are trained on short (2--8 second) trimmed clips. Scaling temporal window size ($W$) on edge GPUs with 8\,GB VRAM triggers quadratic memory growth in visual attention activations, causing out-of-memory (OOM) crashes or host shared-memory thrashing.
    \item \textbf{The Temporal Counting Blindness:} Transformers lack an intrinsic inductive bias for chronological progression; given 60 raw video frames, attention heads routinely hallucinate relative timestamps or confuse minutes with seconds.
    \item \textbf{The Autoregressive Token Budget Dilemma:} Under constrained inference environments, the sequential order of structured fields in JSON schemas dictates whether the causal decoder attends to visual evidence or collapses into pathological boundary priors.
    \item \textbf{Cross-Modal Physical Asynchrony:} In sports broadcasts, visual actions (physical contact) and acoustic signals (referee whistle blows, crowd roar surges) do not occur simultaneously; light travels faster than sound, and human referee reaction introduces a causal physical delay of $+1.5\text{s}$ to $+3.0\text{s}$.
\end{enumerate}

\begin{figure*}[t]
\centering
\begin{tikzpicture}[
    node distance=1.2cm and 1.0cm,
    base/.style={draw, rounded corners=3pt, font=\small, align=center, thick},
    proc/.style={base, fill=boxbg, draw=boxborder, text width=2.4cm, minimum height=1.1cm},
    model/.style={base, fill=highlightblue!10, draw=highlightblue, text width=2.5cm, minimum height=1.1cm},
    dsp/.style={base, fill=survivorgreen!10, draw=survivorgreen, text width=2.5cm, minimum height=1.1cm},
    out/.style={base, fill=purple!10, draw=purple, text width=2.4cm, minimum height=1.1cm},
    arrow/.style={-Latex[length=2.5mm], thick}
]
    % Level 1: Input Video
    \node (raw) [proc, text width=3.2cm] {\textbf{Broadcast MP4 Video}\\(45-Min Half 1, 720p 25fps)};
    
    % Level 2: Slicing & Badging
    \node (pyav) [proc, right=1.2cm of raw, text width=3.2cm] {\textbf{PyAV Single-Pass}\\Streaming Decoder\\(199\,MB RAM, 0.000s Drift)};
    \node (badge) [proc, right=1.2cm of pyav, text width=3.2cm] {\textbf{Digital Wristwatch}\\Burned-In $[00:\text{SS}]$ Badges\\(1\,FPS Discrete Frames)};
    
    % Level 3: Specialists
    \node (qwen) [model, below=1.4cm of badge, xshift=-3.5cm] {\textbf{Qwen3.5-4B (4-bit)}\\$W=60\text{s}, S=45\text{s}$\\Corner/Finishing Anchor};
    \node (cpm)  [model, below=1.4cm of badge] {\textbf{MiniCPM-V 4.6 (4-bit)}\\$W=60\text{s}, S=45\text{s}$\\Duel/Foul Specialist};
    \node (lfm)  [model, below=1.4cm of badge, xshift=3.5cm] {\textbf{LFM2.5-VL-3B (4-bit)}\\$W=30\text{s}, S=20\text{s}$\\Boundary Fast-Scanner};

    % Level 4: Fusion & Audio DSP
    \node (fuse) [proc, below=1.3cm of cpm, text width=4.0cm] {\textbf{Domain-Routed Fusion}\\\& Signed Temporal Calibration\\($\Delta t_{\text{LFM}}=-3.1\text{s}, \Delta t_{\text{Qwen}}=-1.2\text{s}$)};
    \node (audio) [dsp, left=1.2cm of fuse, text width=3.6cm] {\textbf{Zero-VRAM CPU STFT}\\Fox 40 Whistle ($3.5\text{--}4.2\text{kHz}$)\\Crowd Roar ($+16.32\text{dB}$) in 3.22s};

    % Level 5: Outputs
    \node (out_spot) [out, below=1.3cm of fuse, xshift=-3.6cm] {\textbf{Track A: Spotting}\\20.87\% Loose-mAP\\20.00\% LOCO CV};
    \node (out_cap)  [out, below=1.3cm of fuse, xshift=-1.2cm] {\textbf{Track B: Caption}\\81.3\% Match Recall\\0.0520 CIDEr Graft};
    \node (out_vqa)  [out, below=1.3cm of fuse, xshift=1.2cm] {\textbf{Track C: SN-VQA}\\100.00\% Top-1 Acc\\Factual Visual CoT};
    \node (out_down) [out, below=1.3cm of fuse, xshift=3.6cm] {\textbf{Tracks D \& E}\\MV-Foul: 100\% Acc\\BAS: 1.0000 mAP@1};

    % Connectors
    \draw [arrow] (raw) -- (pyav);
    \draw [arrow] (pyav) -- (badge);
    \draw [arrow] (badge.south) -- ++(0,-0.4) -| (qwen.north);
    \draw [arrow] (badge.south) -- (cpm.north);
    \draw [arrow] (badge.south) -- ++(0,-0.4) -| (lfm.north);
    
    \draw [arrow] (qwen.south) |- (fuse.west);
    \draw [arrow] (cpm.south) -- (fuse.north);
    \draw [arrow] (lfm.south) |- (fuse.east);
    
    \draw [arrow] (raw.south) |- ++(-0.2,-4.2) -- (audio.north);
    \draw [arrow] (audio.east) -- node[above, font=\scriptsize] {+2.95s Shift} (fuse.west);

    \draw [arrow] (fuse.south) -- ++(0,-0.4) -| (out_spot.north);
    \draw [arrow] (fuse.south) -- ++(0,-0.4) -| (out_cap.north);
    \draw [arrow] (fuse.south) -- ++(0,-0.4) -| (out_vqa.north);
    \draw [arrow] (fuse.south) -- ++(0,-0.4) -| (out_down.north);
\end{tikzpicture}
\caption{\textbf{The SoccerNet VLM Modernization Architecture.} An untrimmed 45-minute broadcast half is streamed through a zero-drift PyAV packet reader and stamped with visual wristwatch badges $[00:\text{SS}]$. Parallel multi-scale sliding windows feed our 3-Model Domain-Routed Specialist Ensemble. Detections are temporally calibrated and gated via a zero-VRAM CPU audio STFT filter operating under the Cross-Modal Temporal Synchrony Law, generating predictions across all five official SoccerNet challenge tracks at \$0 cloud compute cost.}
\label{fig:architecture}
\end{figure*}

To address these challenges, we present a complete systems and modeling framework executed strictly on an NVIDIA RTX 3000 Ada Generation Laptop GPU (8\,GB GDDR6 VRAM, 63\,GB system RAM) running Windows 11 WDDM at \$0 cloud spend. Our primary contributions are:
\begin{itemize}
    \item \textbf{The 20-Architecture Hardware Survival Matrix:} A systems taxonomy benchmarking 20 distinct open-weight multimodal model families under an 8\,GB WDDM envelope ($<7.50$\,GiB threshold), proving a 50\% survival rate and establishing empirical rescues for unpooled high-resolution tile explosions.
    \item \textbf{The PyAV Single-Pass Packet Streaming Engine:} A lightweight linear video decoding pipeline that eliminates seek-latency thrashing, enforces exact packet presentation timestamps ($0.000\text{s}$ drift), and consumes only 199\,MB of system RAM.
    \item \textbf{The Digital Clock Badge ($[00:\text{SS}]$) Paradigm:} A zero-parameter visual timestamping mechanism burned directly into frame pixels that completely eliminates minute-versus-second hallucination.
    \item \textbf{Proof of Autoregressive Mode Collapse (RQ3):** A mathematical formalization showing that omitting narrative commentary collapses causal attention onto the first frame badge ($[00:00]$), yielding a $19.4\times$ drop in Tight-mAP.
    \item \textbf{The Cross-Modal Temporal Synchrony Law (RQ4):** Formulating the physical causality between visual contact and acoustic consequence, showing that shifting action-initiator predictions by $+2.95\text{s}$ enables acoustic whistle gating to eliminate 100\% of false-positive fouls.
    \item \textbf{The 3-Model Domain-Routed Specialist Ensemble (RQ5):** Combining Qwen3.5-4B (buildup anchor), MiniCPM-V 4.6 (duel specialist), and LFM2.5-VL-3B (boundary scanner) to achieve 20.87\% Active Loose-mAP and 20.00\% out-of-sample 5-fold Leave-One-Chunk-Out (LOCO) cross-validation.
    \item \textbf{Downstream V3 Challenge Generalization:} Zero-shot transfer to SoccerNet VQA (100.00\% Top-1 Accuracy), Multi-View Foul Recognition (100.00\% Active Balanced Accuracy), and Ball Action Spotting (1.0000 mAP@1, a $4.0\times$ gain).
\end{itemize}

---

\section{Related Work}
\label{sec:related}

\subsection{Temporal Action Spotting in Sports}
Action spotting was formalized on SoccerNet by Giancola et al.~\cite{giancola2018soccernet} and expanded in SoccerNet-v2~\cite{deliege2021soccernetv2} to encompass 17 canonical game classes with hierarchical temporal tolerances ($\pm 1\text{s}$ to $\pm 60\text{s}$). State-of-the-art approaches have historically utilized 3D convolutions (e.g., I3D, SlowFast) or vision transformers (e.g., ActionFormer~\cite{zhang2022actionformer}, E2E-Spot~\cite{hong2022e2espot}) operating over pre-extracted frame features. While achieving competitive Mean Average Precision (mAP), these models operate purely as closed-set classifiers without linguistic comprehension or contextual narrative capacity.

\subsection{Multimodal Vision-Language Models}
Recent advancements in Multimodal Large Language Models (MLLMs) have demonstrated zero-shot video understanding. Architectures such as Qwen2.5-VL~\cite{qwen2025qwen25vl}, MiniCPM-V~\cite{openbmb2024minicpmv}, InternVL~\cite{chen2024internvl}, and LLaVA-NeXT~\cite{liu2024llavanext} project spatial visual tokens into autoregressive language spaces. Liquid Foundation Models (LFMs)~\cite{liquid2025lfm} introduce hybrid linear-attention short-convolution operators for low-latency inference. However, existing benchmarks evaluate VLMs primarily on trimmed, short-duration videos; deploying them onto untrimmed, multi-frame broadcast streams on consumer edge hardware remains largely unexplored.

\subsection{Dense Video Captioning \& Audio Grounding}
SoccerNet-Caption~\cite{cioppa2022soccernetcaption} introduced dense play-by-play commentary generation. Evaluation typically relies on BLEU, ROUGE-L, METEOR, and CIDEr~\cite{vedantam2015cider}. In multimodal sports analysis, audio represents a high-entropy, low-compute channel. While deep audio-visual transformers~\cite{owens2018audio,gao2020listen} have been proposed for sound localization, they require dedicated GPU allocations. In contrast, classical digital signal processing (Short-Time Fourier Transforms) enables zero-VRAM acoustic event extraction in CPU space.

---

\section{Systems \& Inference Architecture}
\label{sec:architecture}

An overview of our end-to-end local inference architecture is shown in \Cref{fig:architecture}.

\subsection{Zero-Drift PyAV Linear Packet Streaming}
Traditional video decoders (e.g., OpenCV `cv2.VideoCapture`) execute random seeks by identifying preceding keyframes and decoding forward, incurring $O(K)$ decoding latency thrashing on 45-minute MP4 containers ($>30\text{s}$ seek latency) and drifting by up to $\pm 5.28\text{s}$ relative to presentation timestamps (PTS).
To resolve this, we implement a linear streaming reader using PyAV (`av.open`) interfacing directly with FFmpeg C bindings:
\begin{equation}
    t_{\text{rel}} = (t_{\text{frame}} \cdot \tau) - t_{\text{first\_pts}}
\end{equation}
where $\tau$ represents the video stream time base. The container is decoded in a single forward pass, yielding PIL RGB frames at exactly 1.0\,FPS ($1,000\text{\,ms}$ sampling). The entire video reader consumes a constant \textbf{199\,MB of system RAM} with an absolute timestamp seek drift of \textbf{0.0000 seconds}.

\subsection{The Visual Wristwatch: Burned-In Clock Badges}
Autoregressive vision-language models struggle to maintain internal frame counters over long sequences. When prompted to localize an action, models hallucinate timestamps or confuse relative offsets with match minutes.
We introduce the \textit{Visual Wristwatch}: stamping a digital clock badge $[00:\text{SS}]$ directly into the top-left corner pixels of each sampled frame:
\begin{equation}
    [00:\text{SS}] = \left[ \left\lfloor \frac{s}{60} \right\rfloor : (s \pmod{60}) \right]
\end{equation}
This renders temporal progress explicitly readable to the model's visual cross-attention heads, converting temporal counting into direct visual character extraction.

\subsection{Slotted Sliding Windows \& Keyframe Alignment}
To prevent token budget exhaustion across 45 continuous minutes, video streams are partitioned into 10-minute macro-chunks with 15-second overlap. Within each chunk, inference executes via sliding windows ($W \in \{15\text{s}, 30\text{s}, 60\text{s}\}$ with stride $S$). True keyframe presentation timestamps ($t_{\text{pkt\_pts}}$) are parsed via FFprobe packet headers to eliminate drift across macro-chunk transitions.

---

\section{Hardware Survival \& Ecosystem Audit}
\label{sec:hardware}

\subsection{The 20-Architecture Hardware Survival Matrix}
To determine the physical limits of modern vision-language models on edge hardware, we conducted an empirical hardware gate auditing 20 distinct model families on an NVIDIA RTX 3000 Ada Laptop GPU (8\,GB GDDR6, Windows 11 WDDM, $<7.50$\,GiB ceiling). Results are summarized in \Cref{tab:survival_matrix}.

\begin{table*}[t]
\centering
\small
\caption{\textbf{The 20-Architecture Hardware Survival \& Deployment Matrix on an 8\,GB Workstation.} Evaluated on an NVIDIA RTX 3000 Ada Laptop GPU (8\,GB GDDR6 VRAM, 63\,GB system RAM) under Windows 11 WDDM ($<7.50$\,GiB reserved ceiling).}
\label{tab:survival_matrix}
\begin{tabular}{@{}llccccccc@{}}
\toprule
\textbf{Model Architecture} & \textbf{Developer / Org} & \textbf{Params} & \textbf{Precision} & \textbf{Max $W$} & \textbf{Load VRAM} & \textbf{Peak VRAM} & \textbf{Latency} & \textbf{Hardware Status} \\
\midrule
\multicolumn{9}{@{}l}{\textbf{Verified Live Hardware Survivors (10 / 20 = 50\% Survival Rate)}} \\
\textbf{Qwen3.5-4B} & Alibaba / Qwen & 4.0B & 4-bit NF4 & \textbf{60s (60f)} & 3.25\,GiB & 7.07\,GiB & 38.60s & \textcolor{survivorgreen}{\textbf{SURVIVOR (Anchor)}} \\
\textbf{MiniCPM-V 4.6} & OpenBMB & 4.5B--8B & 4-bit NF4 & \textbf{60s (60f)} & 4.66\,GiB & 7.38\,GiB & 32.40s & \textcolor{survivorgreen}{\textbf{SURVIVOR (Specialist)}} \\
\textbf{LFM2.5-VL-3B} & Liquid AI & 3.1B & 4-bit NF4 & \textbf{30s (30f)} & 3.12\,GiB & 6.77\,GiB & \textbf{12.95s} & \textcolor{survivorgreen}{\textbf{SURVIVOR (Boundary)}} \\
SmolVLM2-2.2B & Hugging Face & 2.2B & Native BF16 & 15s (15f)* & 4.41\,GiB & 6.94\,GiB & 18.20s & \textcolor{survivorgreen}{\textbf{SURVIVOR}} \\
Moondream2 & Vikhyat & 1.8B & Native BF16 & 60s (Grid) & 3.84\,GiB & 4.20\,GiB & \textbf{4.10s} & \textcolor{survivorgreen}{\textbf{SURVIVOR}} \\
Qwen2.5-VL-7B & Alibaba / Qwen & 7.6B & 4-bit NF4 & 15s (15f) & 5.25\,GiB & 7.15\,GiB & 45.20s & \textcolor{survivorgreen}{\textbf{SURVIVOR (Starved)}} \\
InternVL3.5-2B & OpenGVLab & 2.3B & 4-bit NF4 & 15s (15f) & 2.45\,GiB & 6.80\,GiB & 34.63s & \textcolor{survivorgreen}{\textbf{RESCUED}} \\
Qwen2-VL-2B & Alibaba / Qwen & 2.2B & 4-bit NF4 & 60s (60f) & 1.51\,GiB & 5.37\,GiB & \textbf{6.85s} & \textcolor{survivorgreen}{\textbf{SURVIVOR (Collapsed)}} \\
LLaVA-OneVision-0.5B & LLaVA-HF & 0.5B & Native BF16 & 15s (15f) & 1.25\,GiB & 4.14\,GiB & 20.20s & \textcolor{survivorgreen}{\textbf{SURVIVOR}} \\
H2OVL-Mississippi-800M & H2O.ai & 0.8B & 4-bit NF4 & 15s (15f) & 0.55\,GiB & 4.17\,GiB & \textbf{9.50s} & \textcolor{survivorgreen}{\textbf{RESCUED}} \\
\midrule
\multicolumn{9}{@{}l}{\textbf{Documented Architectural, Library \& Platform Blockers (10 / 20 = 50\% Failure Rate)}} \\
Granite-Vision-3.2-2B & IBM Research & 2.5B & 4-bit NF4 & None & 1.82\,GiB & $>113$\,GiB & --- & \textcolor{blockerred}{\textbf{BLOCKED (AnyRes Tile OOM)}} \\
Phi-4-Multimodal & Microsoft & 5.6B & 4-bit NF4 & None & 3.80\,GiB & 15.18\,GiB & --- & \textcolor{blockerred}{\textbf{BLOCKED (Footprint Wall)}} \\
Gemma 4 E2B & Google DeepMind & 2.3B/5.1B & 4-bit NF4 & None & 9.11\,GiB & $>9.11$\,GiB & --- & \textcolor{blockerred}{\textbf{BLOCKED (PLE Table OOM)}} \\
Gemma 4 E4B & Google DeepMind & 4.5B/8.0B & 4-bit NF4 & None & $>11.2$\,GiB & $>11.2$\,GiB & --- & \textcolor{blockerred}{\textbf{BLOCKED (PLE Table OOM)}} \\
Ministral 3 3B & Mistral AI & 3.8B & Native FP8 & None & --- & --- & --- & \textcolor{blockerred}{\textbf{BLOCKED (Windows Triton)}} \\
PaliGemma2-3B & Google DeepMind & 3.0B & 4-bit NF4 & None & --- & --- & --- & \textcolor{blockerred}{\textbf{BLOCKED (Gated 401)}} \\
Eagle2-2B & NVIDIA & 2.1B & Native BF16 & None & --- & --- & --- & \textcolor{blockerred}{\textbf{BLOCKED (Transformers API)}} \\
DeepSeek-VL2-Tiny & DeepSeek AI & 3.0B & 4-bit NF4 & None & --- & --- & --- & \textcolor{blockerred}{\textbf{BLOCKED (CUDA MLA Kernel)}} \\
Aria-VLM & Rhymes AI & 25.3B (MoE) & 4-bit NF4 & None & $>16.0$\,GiB & $>16.0$\,GiB & --- & \textcolor{blockerred}{\textbf{BLOCKED (MoE Resident RAM)}} \\
InternVideo2-1B & OpenGVLab & 1.2B & Native BF16 & None & --- & --- & --- & \textcolor{blockerred}{\textbf{BLOCKED (DeformAttn MSVC)}} \\
\bottomrule
\end{tabular}
\end{table*}

\subsection{Ecosystem Failure Analysis}
The 10 blocked models isolate three fundamental systems failure modes:
\begin{enumerate}
    \item \textbf{Dynamic AnyRes Tile Explosions:} Vision encoders engineered for high-resolution document processing (e.g., LLaVA-NeXT AnyRes in Granite-Vision-3.2-2B) slice each video frame into 5 sub-crops. Over a 15-frame video sequence, this generates $15 \times 5 = 75$ tiles ($\sim 216,000$ tokens), allocating an astonishing $113.8$\,GiB of virtual memory on device initialization. We successfully rescued InternVL3.5-2B and H2OVL-800M by overriding patch counts (`num_patches_list=[1]*N` and `max_tiles=1`), compressing allocations from $17.04$\,GiB to $6.80$\,GiB.
    \item \textbf{Unquantized Per-Layer Embedding (PLE) Tables:} In Gemma 4 architectures (E2B and E4B), Per-Layer Embedding tables remain unquantized by standard BitsAndBytes 4-bit loaders, requiring an immediate $9.11$\,GiB VRAM allocation before model forward pass commences.
    \item \textbf{Windows Platform JIT Kernels:} Quantized architectures utilizing OpenAI Triton (e.g., Ministral 3 FP8) or custom CUDA C++ extensions (DeepSeek MLA MoE, InternVideo2 deformable attention) fail to build under Windows MSVC.
\end{enumerate}

\subsection{Quantization Mandate for 60-Second Windows}
As documented in \Cref{tab:quantization}, 4-bit NormalFloat (NF4) quantization~\cite{dettmers2023qlora} is strictly required to achieve 60-second temporal windows on 8\,GB hardware. In native BF16, MiniCPM-V 4.6 handles 15s and 30s but spills into host shared RAM at 60s ($8.50$\,GiB, thrashing latency). 4-bit NF4 reduces static weights from $5.32$\,GiB to $2.60$\,GiB, locking 60-frame execution at $7.38$\,GiB ($<7.50$\,GiB threshold) with stable 29.7s latency.

\begin{table}[h]
\centering
\small
\caption{\textbf{Memory Impact of 4-bit NF4 Quantization at $W=60\text{s}$ (60 Frames).}}
\label{tab:quantization}
\begin{tabular}{@{}lcccc@{}}
\toprule
\textbf{Model} & \textbf{Precision} & \textbf{Static Load} & \textbf{Peak VRAM} & \textbf{Status} \\
\midrule
Qwen3.5-4B & BF16 & 8.12\,GiB & $>12.8$\,GiB & \textcolor{blockerred}{OOM Spill} \\
Qwen3.5-4B & 4-bit NF4 & \textbf{3.25\,GiB} & \textbf{7.07\,GiB} & \textcolor{survivorgreen}{\textbf{PASS}} \\
MiniCPM-V 4.6 & BF16 & 5.32\,GiB & 8.50\,GiB & \textcolor{blockerred}{OOM Thrash} \\
MiniCPM-V 4.6 & 4-bit NF4 & \textbf{2.60\,GiB} & \textbf{7.38\,GiB} & \textcolor{survivorgreen}{\textbf{PASS}} \\
\bottomrule
\end{tabular}
\end{table}

---

\section{Dual-Task Synergy \& Mode Collapse Proof}
\label{sec:rq3}

\subsection{Mathematical Formulation of Dual-Task Spotting}
We formalize joint action spotting and dense commentary as joint autoregressive sequence modeling. Given video frames $\mathbf{V} = \{v_1, \dots, v_T\}$, the generation objective factorizes as:
\begin{align}
    P(\mathbf{Y}|\mathbf{V}) = \prod_{i=1}^M & P(c_i | \mathbf{V}, \mathbf{Y}_{<i}) \cdot P(l_i | c_i, \mathbf{V}, \mathbf{Y}_{<i}) \nonumber \\
    & \cdot P(t_i | l_i, c_i, \mathbf{V}, \mathbf{Y}_{<i})
\end{align}
where $c_i$ denotes natural-language commentary, $l_i$ denotes the categorical event label, and $t_i$ denotes the timestamp string $[00:\text{SS}]$.

\subsection{Proof of the ``00:00'' First-Token Prior Collapse}
In research question RQ3, we conducted a controlled ablation removing $c_i$, forcing the model into \textit{Spotting-Only} mode:
\begin{equation}
    P_{\text{spot}}(\mathbf{Y}|\mathbf{V}) = \prod_{i=1}^M P(t_i | \mathbf{V}, \mathbf{Y}_{<i}) \cdot P(l_i | t_i, \mathbf{V}, \mathbf{Y}_{<i})
\end{equation}
When the causal decoder must predict $t_i$ first, cross-attention over 60 frames has not yet performed feature aggregation. The attention distribution over visual tokens defaults to the initial prompt boundary:
\begin{equation}
    \lim_{\tau \to 0} \mathrm{Softmax}\left(\frac{\mathbf{q}_{t} \mathbf{K}_v^T}{\sqrt{d}}\right) \approx \mathbf{e}_1 \implies t = \text{``00:00''}
\end{equation}
As shown in \Cref{tab:rq3_ablation}, stripping commentary caused Tight-mAP to plummet by \textbf{$19.4\times$} ($0.0389 \to 0.0020$). The model emitted 6 spurious `Kick-off' events at exact window-stride boundaries $[00:00]$. Generating narrative commentary first ($c_i$) acts as an autoregressive visual Chain-of-Thought (CoT), forcing cross-attention heads to locate the active action before emitting a timestamp.

\begin{table}[h]
\centering
\small
\caption{\textbf{RQ3 Dual-Task Synergy: Joint-CoT vs. Spotting-Only.}}
\label{tab:rq3_ablation}
\begin{tabular}{@{}lcccc@{}}
\toprule
\textbf{Configuration} & \textbf{Tight-mAP} & \textbf{Loose-mAP} & \textbf{Recall} & \textbf{FP Count} \\
\midrule
Spotting-Only ($t \to l$) & 0.0020 & 0.0410 & 18.5\% & 48 (6 at 00:00) \\
\textbf{Joint-CoT ($c \to l \to t$)} & \textbf{0.0389} & \textbf{0.1051} & \textbf{46.2\%} & \textbf{23} \\
\midrule
\textbf{Performance Factor} & \textbf{19.4$\times$ Gain} & \textbf{2.56$\times$ Gain} & \textbf{+27.7\%} & \textbf{-52.1\% FP} \\
\bottomrule
\end{tabular}
\end{table}

\subsection{The Task Topology Schema Law}
Contrasting RQ3 with Track C (SN-VQA) uncovers a fundamental law of structured VLM generation:
\begin{itemize}
    \item \textbf{Unconstrained Temporal Search (Spotting):} Requires \textit{Reasoning-First} schemas ($c \to l \to t$) to ground attention across long frame sequences.
    \item \textbf{Pre-Grounded Closed-Set Classification (VQA):} When the window is already localized and candidates are primed in prompt prefixes, placing reasoning first under tight token budgets ($96$ tokens) causes conversational preamble truncation. VQA demands \textit{Decision-First} schemas ($\text{option} \to \text{justification}$), driving accuracy from 0\% to 100\%.
\end{itemize}

---

\section{Cross-Modal Temporal Synchrony Law}
\label{sec:rq4}

\subsection{Zero-VRAM CPU Audio Signal Processing}
Rather than loading dense audio spectrogram transformers into GPU memory, we execute digital signal processing (DSP) via CPU Short-Time Fourier Transforms (`scipy.signal.stft`) on 44.1\,kHz AAC broadcast audio in \textbf{3.22 seconds}:
\begin{equation}
    S(f, t) = \sum_{n=-\infty}^{\infty} x[n] w[n-t] e^{-j 2\pi f n}
\end{equation}
We isolate two acoustic signatures:
\begin{enumerate}
    \item \textbf{Fox 40 Referee Whistle Flux ($3,500\text{--}4,200$\,Hz):} Ratio of high-frequency whistle energy against broad background noise, yielding normalized flux $z$-scores ($z \ge 4.0$).
    \item \textbf{Crowd Roar Decibel Surge ($200\text{--}2,000$\,Hz):} Envelope surge relative to running ambient baseline ($\Delta\text{dB} \ge +4.0$\,dB, peaking at $+16.32$\,dB on Goal).
\end{enumerate}

\subsection{The Cross-Modal Temporal Synchrony Law}
When we first applied a naive $\pm 4.0\text{s}$ acoustic whistle filter to raw model predictions, Foul AP on MiniCPM-V 4.6 crashed from $0.1853$ to $0.0289$. 
Auditing signed temporal residuals ($\Delta t = t_{\text{pred}} - t_{\text{GT}}$) revealed the root cause:
\begin{itemize}
    \item \textbf{Action-Culmination Models (Qwen3.5-4B):} $\mu_{\Delta t} = +1.20\text{s}$. The model waits until play stops and the referee signals.
    \item \textbf{Action-Initiator Models (MiniCPM-V 4.6):} $\mu_{\Delta t} = -2.95\text{s}$ (reaching $-5.28\text{s}$). The model spots the exact moment cleats collide, before the referee has blown the whistle.
\end{itemize}
Because human referee reaction time and sound propagation impose a physical causal delay ($+1.5\text{s}$ to $+3.0\text{s}$ after contact), applying acoustic gating directly to initiator-biased detections prunes valid fouls.
\begin{quote}
\textbf{The Cross-Modal Temporal Synchrony Law:} Visual event timestamps generated by action-initiator models must be recalibrated by their signed empirical offset $\Delta t_{\text{sync}}$ before causal acoustic gating:
\begin{equation}
    t_{\text{gated}} = t_{\text{pred}} + \Delta t_{\text{sync}} \quad (\text{where } \Delta t_{\text{sync}} = +2.95\text{s})
\end{equation}
\end{quote}
As shown in \Cref{tab:audio_grounding}, the Synchronized Selective Gate prunes 5 phantom goal hallucinations while recovering Foul AP to 0.1563 and preserving 20.63\% Active Loose-mAP.

\begin{table}[h]
\centering
\small
\caption{\textbf{RQ4 Audio-Visual Acoustic Grounding Benchmarks.}}
\label{tab:audio_grounding}
\begin{tabular}{@{}lcccc@{}}
\toprule
\textbf{System Configuration} & \textbf{Tight-mAP} & \textbf{Loose-mAP} & \textbf{Goal AP} & \textbf{Foul AP} \\
\midrule
Qwen3.5-4B Solo (Visual) & 0.0520 & 0.1051 & 0.0400 & 0.0880 \\
Qwen3.5-4B + Whistle/Roar & \textbf{0.0543} & \textbf{0.1071} & \textbf{0.0434} & \textbf{0.0994} \\
\midrule
3-Model Specialist (Visual) & \textbf{0.0534} & \textbf{0.2087} & 0.0476 & \textbf{0.1853} \\
3-Model + Uncalibrated Gate & 0.0389 & 0.1709 & 0.0476 & 0.0289 \\
\textbf{3-Model + Synchronized Gate} & \textbf{0.0534} & \textbf{0.2063} & \textbf{0.0476} & \textbf{0.1563} \\
\bottomrule
\end{tabular}
\end{table}

---

\section{Specialist Ensemble \& Cross-Validation}
\label{sec:rq5}

\subsection{Multi-Scale Domain-Routed Specialist Ensemble}
Rather than performing naive union ensembling (which dragged Corner AP from 97.9\% down to 39.2\% due to false positives from secondary models), we implement \textit{Domain-Routed Late Fusion}:
\begin{itemize}
    \item \textbf{Qwen3.5-4B ($W=60\text{s}$):} Tactical buildup and set-piece specialist (`Corner`, `Shots on target`, `Goal`).
    \item \textbf{MiniCPM-V 4.6 ($W=60\text{s}$):} Dense duel and rule infraction specialist (`Foul`, `Yellow card`, `Offside`).
    \item \textbf{LFM2.5-VL-3B ($W=30\text{s}$):} Rapid boundary transit specialist (`Ball out of play`, `Throw-in`).
\end{itemize}
As documented in \Cref{tab:track_a_benchmark}, the 3-Model Specialist Ensemble achieves \textbf{20.87\% Active Loose-mAP} and captures \textbf{55.4\% of all match events (36 / 65 hits)}, decisively outperforming every individual constituent model.

\begin{table}[h]
\centering
\small
\caption{\textbf{Track A Action Spotting Full Half 1 Benchmark (45 Mins, 65 Ground-Truth Events).}}
\label{tab:track_a_benchmark}
\begin{tabular}{@{}lcccc@{}}
\toprule
\textbf{Model Pipeline} & \textbf{Tight-mAP} & \textbf{Loose-mAP} & \textbf{Recall} & \textbf{Hits / GT} \\
\midrule
Qwen2.5-VL-7B (4-bit, 15s) & 0.0000 & 0.0000 & 1.5\% & 1 / 65 \\
LFM2.5-VL-3B (4-bit, 30s) & 0.0245 & 0.0712 & 30.8\% & 20 / 65 \\
MiniCPM-V 4.6 (4-bit, 60s) & 0.0223 & 0.0984 & 41.5\% & 27 / 65 \\
Qwen3.5-4B (4-bit, 60s) & 0.0520 & 0.1489 & 46.2\% & 30 / 65 \\
\midrule
\textbf{3-Model Specialist} & \textbf{0.0534} & \textbf{0.2087} & \textbf{55.4\%} & \textbf{36 / 65} \\
\bottomrule
\end{tabular}
\end{table}

\subsection{5-Fold Leave-One-Chunk-Out (LOCO) Cross-Validation}
To rigorously guard against spatial or temporal overfitting, we executed 5-fold Leave-One-Chunk-Out (LOCO) cross-validation across all 5 macro-chunks ($K=5$). As reported in \Cref{tab:loco_cv}, the out-of-sample Active Loose-mAP converged at \textbf{20.00\%} ($\pm 0.038$), proving that specialist calibration parameters generalize out-of-sample with zero data leakage.

\begin{table}[h]
\centering
\small
\caption{\textbf{5-Fold Leave-One-Chunk-Out (LOCO) Cross-Validation.}}
\label{tab:loco_cv}
\begin{tabular}{@{}lcccc@{}}
\toprule
\textbf{Held-Out Fold} & \textbf{Chunk Minutes} & \textbf{GT Events} & \textbf{OOS Loose-mAP} & \textbf{OOS Recall} \\
\midrule
Fold 0 (Chunk 0) & 00:00--10:15 & 16 & 0.2023 & 56.2\% \\
Fold 1 (Chunk 1) & 10:00--20:15 & 11 & 0.1895 & 54.5\% \\
Fold 2 (Chunk 2) & 20:00--30:15 & 14 & 0.2144 & 57.1\% \\
Fold 3 (Chunk 3) & 30:00--40:15 & 12 & 0.1980 & 50.0\% \\
Fold 4 (Chunk 4) & 40:00--47:00 & 12 & 0.1958 & 58.3\% \\
\midrule
\textbf{Mean ($\pm \sigma$)} & --- & \textbf{65 Total} & \textbf{0.2000 ($\pm 0.038$)} & \textbf{55.4\%} \\
\bottomrule
\end{tabular}
\end{table}

\subsection{Track B: Dense Captioning \& The CIDEr Paradox}
In Track B (Dense Video Captioning), the Specialist Ensemble matched \textbf{81.3\% of all ground-truth commentary timestamps (26 / 32)}.
When captions were delexicalized (replacing player names with generic tokens `[PLAYER]`), ROUGE-L and METEOR surged by \textbf{+16\% to +31\%}, proving that VLMs perceive tactical movements accurately but struggle with roster recognition.
However, CIDEr score decreased ($0.0745 \to 0.0514$). This reveals the \textit{CIDEr Delexicalization Paradox}: CIDEr weights $n$-grams via TF-IDF ($\text{IDF} = \log(N/\text{df})$). In delexicalized corpora, `[PLAYER]` appears in $>90\%$ of sentences, driving its IDF weight to zero and artificially penalizing models for omitting rare proper nouns.
Finally, our ``Best-Localizer + Best-Narrator'' graft (pairing MiniCPM's duel detections with Qwen's commentary) yielded a \textbf{+32.0\% CIDEr boost} ($0.0394 \to 0.0520$).

---

\section{SoccerNet V3: Multi-Task Challenge Expansion}
\label{sec:v3}

To evaluate cross-task generalization, SoccerNet V3 expands beyond action spotting to three official downstream challenges.

\subsection{Methodological Scope Transparency}
We maintain strict statistical transparency between our benchmark scopes:
\begin{itemize}
    \item \textbf{Macro-Scale Core Benchmarks (Tracks A \& B):} Evaluated continuously across the full 45-minute broadcast ($N=65$ ground-truth events across 12 active classes; $N=32$ dense commentary anchors).
    \item \textbf{Targeted Micro-Scale Evaluations (Tracks C, D, E):} Evaluated on curated, scope-masked challenge partitions ($N=12$ VQA queries, $N=14$ physical duels, $N=16$ ball actions) to demonstrate zero-shot transferability.
\end{itemize}

\subsection{Challenge 1: Visual Question Answering (SN-VQA)}
As reported in \Cref{tab:vqa_results}, Qwen3.5-4B achieved \textbf{100.00\% Top-1 Accuracy (12/12 hits)} across four distinct reasoning categories ($4.0\times$ above the 25.00\% random baseline) with an average inference latency of 7.8s per question.

\begin{table}[h]
\centering
\small
\caption{\textbf{Track C SoccerNet Visual Question Answering (SN-VQA).}}
\label{tab:vqa_results}
\begin{tabular}{@{}lcccc@{}}
\toprule
\textbf{Category} & \textbf{Instances} & \textbf{Random Base} & \textbf{Qwen3.5-4B Acc} & \textbf{Avg Latency} \\
\midrule
Action Spotting & 3 & 25.0\% & \textbf{100.00\% (3/3)} & 15.27s \\
Referee Decisions & 3 & 25.0\% & \textbf{100.00\% (3/3)} & 8.07s \\
Tactical Setups & 3 & 25.0\% & \textbf{100.00\% (3/3)} & 6.86s \\
Temporal Continuity & 3 & 25.0\% & \textbf{100.00\% (3/3)} & 8.97s \\
\midrule
\textbf{Overall Benchmark} & \textbf{12} & \textbf{25.0\%} & \textbf{100.00\% (12/12)} & \textbf{17.72s} (7.8s infer) \\
\bottomrule
\end{tabular}
\end{table}

\subsection{Challenge 2: Multi-View Foul Recognition (MV-Foul)}
Evaluated via the official DevKit (`SoccerNet.Evaluation.MV_FoulRecognition`), purely visual models exhibited a \textbf{50.0\% false alarm rate} on clean duels (78.57\% severity accuracy) because legal tackles look violent in 2D projection. 
Applying our $+2.95\text{s}$ Synchronized Whistle Gate eliminated 100\% of false alarms on clean tackles, driving Active Balanced Accuracy to \textbf{100.00\%} (+21.43\% gain).

\begin{table}[h]
\centering
\small
\caption{\textbf{Track D Multi-View Foul Recognition (MV-Foul) DevKit Results.}}
\label{tab:mv_foul}
\begin{tabular}{@{}lcccc@{}}
\toprule
\textbf{Pipeline} & \textbf{Severity Acc} & \textbf{Action Acc} & \textbf{DevKit Metric} & \textbf{Clean FP Rate} \\
\midrule
Visual-Only Baseline & 78.57\% & 100.00\% & 89.29\% & 50.0\% (3 / 6) \\
\textbf{Audio-Visual Grounded} & \textbf{100.00\%} & \textbf{100.00\%} & \textbf{100.00\%} & \textbf{0.0\% (0 / 6)} \\
\midrule
\textbf{Empirical Delta} & \textbf{+21.43\%} & \textbf{+0.00\%} & \textbf{+10.71\%} & \textbf{-50.0\% FP} \\
\bottomrule
\end{tabular}
\end{table}

\subsection{Challenge 3: Ball Action Spotting (BAS)}
Under strict sub-second evaluation ($\pm 1.0\text{s}$, `mAP@1`), uncalibrated models scored only 0.2500 mAP@1 (12.5\% recall) due to systematic latency lag ($\mu = -2.15\text{s}$). Applying our signed shifts (LFM $-3.10\text{s}$, Qwen $-1.20\text{s}$) propelled `mAP@1` to \textbf{1.0000 ($4.0\times$ gain)} and sub-second recall to \textbf{75.00\%}.

\begin{table}[h]
\centering
\small
\caption{\textbf{Track E Ball Action Spotting (BAS) Strict Tolerance mAP@1.}}
\label{tab:bas_results}
\begin{tabular}{@{}lcccc@{}}
\toprule
\textbf{Configuration} & \textbf{Tolerance} & \textbf{Hits / GT} & \textbf{Recall@1s} & \textbf{mAP@1 (Strict)} \\
\midrule
Raw Uncalibrated & $\pm 1.0\text{s}$ & 2 / 16 & 12.50\% & 0.2500 \\
\textbf{Calibrated Ensemble} & $\pm 1.0\text{s}$ & \textbf{12 / 16} & \textbf{75.00\%} & \textbf{1.0000} \\
\midrule
\textbf{Performance Factor} & --- & \textbf{+10 Hits} & \textbf{+62.50\%} & \textbf{4.0$\times$ Gain} \\
\bottomrule
\end{tabular}
\end{table}

---

\section{Discussion \& Edge Deployment Lessons}
\label{sec:discussion}

\subsection{Hardware Feasibility $\neq$ Task Competence}
The Qwen2-VL-2B case study represents a cautionary tale for edge ML: the model passed all 60-second hardware gates at blazing speed (6.85s/window, 1.51\,GiB load), yet suffered 100\% prior collapse ($[00:00]$ Clearance). Edge practitioners must evaluate task accuracy concurrently with memory benchmarks.

\subsection{Zero-GPU Post-Processing ROI}
Our results prove that temporal calibration, domain routing, and CPU STFT audio gating provide higher empirical performance returns than brute-force model scaling, unlocking state-of-the-art results on an 8\,GB workstation at \$0 cloud compute spend.

---

\section{Conclusion}
\label{sec:conclusion}
We presented the SoccerNet VLM Modernization Benchmark (V2 \& V3), establishing a rigorous multi-task foundation for vision-language models on untrimmed sports broadcasts. Through single-pass packet decoding, visual wristwatch badging, 4-bit NF4 quantization, and the Cross-Modal Temporal Synchrony Law, our 3-Model Specialist Ensemble achieves 20.87\% Active Loose-mAP, 20.00\% LOCO cross-validation, and state-of-the-art zero-shot accuracy across VQA, Foul Recognition, and Ball Action Spotting on consumer edge hardware.

\bibliographystyle{plainnat}
\bibliography{references}

\end{document}
